\begin{figure}[H]
\centering
\begin{tikzpicture}[
  tokena/.style={draw=petroleo, fill=petroleo!26, rounded corners=2pt,
                 minimum width=1.45cm, minimum height=0.55cm, inner sep=2pt},
  tokenb/.style={draw=vinoCIMAT!70, fill=vinoCIMAT!12, rounded corners=2pt,
                 minimum width=1.45cm, minimum height=0.55cm, inner sep=2pt},
  attention/.style={petroleo, line width=0.85pt},
  every node/.style={font=\small}
]
  \node[anchor=east, font=\bfseries] at (-1.15, 1.4) {Registro A};
  \node[tokena] (a1) at (0, 1.4) {Juan};
  \node[tokena] (a2) at (1.8, 1.4) {45};
  \node[tokena] (a3) at (3.6, 1.4) {21/04};
  \node[tokena] (a4) at (5.4, 1.4) {COVID-19};

  \node[anchor=east, font=\bfseries] at (-1.15, -0.4) {Registro B};
  \node[tokenb] (b1) at (0, -0.4) {J. P\'erez};
  \node[tokenb] (b2) at (1.8, -0.4) {46};
  \node[tokenb] (b3) at (3.6, -0.4) {23/05};
  \node[tokenb] (b4) at (5.4, -0.4) {SARS-CoV};

  \foreach \a in {a1,a2,a3,a4}{
    \foreach \b in {b1,b2,b3,b4}{
      \draw[attention] (\a.south) -- (\b.north);
    }
  }
\end{tikzpicture}
\caption{Visualización conceptual del mecanismo de atenci\'on cruzada en el Cross-Encoder. Las conexiones representan la interacción fina entre los tokens de ambos registros procesados; el diagrama es solo ilustrativo y no corresponde a pesos de atenci\'on observados.}
\label{fig:cross_attention}
\end{figure}
